\scanpage{11}
Il est un peu étrange que l'expression obtenue pour $\Pi_1(S_0\mathcal T_{0,I},\mathfrak S_I;x)$ ne dépende pas du choix de $x\in\omega(I)$ -- c'est un groupe $G(I)$, extension de $\mathfrak S_I$ par $\mathbb Z^I$, qui se construit canoniquement, sans référence à la donnée d'un $x\in\omega(I)$. Donc $\mathfrak S_I$ opère sur ce groupe \inserted{extension} (de façon compatible avec son op. sur $\mathbb Z^I$ et sur lui-$=$), bien qu'il n'y ait pas de scindage par quoi décrire cette opération. Bien entendu, la restriction de cette opération à $\mathfrak S_I^+$ n'est autre que celle déduite de ce scindage canonique au dessus de $\mathfrak S_I^+$.

\figure{\begin{tikzpicture}\draw (0,0)--(1,0);\end{tikzpicture}}

Reprenons maintenant le cas d'un graphe MD\deleted{\unsure{G}} $G$ de dim. modulaire $0$, de sorte que $S_0\mathcal T_G$ se déduit du produit
\[
S_0\mathcal T_{G^b}=\prod_{s\in S}S_0\mathcal T_{G_s^b}
=\prod_{s\in S}S_0\mathcal T_{0,\widehat{\vec A}_s}
\]
par l'hom de groupes \deleted{d'op.} d'op. internes
\[
(20\qquad\mathbb Z^{\widehat{\vec A}}\longrightarrow\mathbb Z^{\widehat A}.
\]
On peut d'autre part décrire le produit $S_0\mathcal T_{G^b}$ en termes du torseur
\[
\unsure{(21)}\qquad\varepsilon(G)=\varepsilon(G^b)=\prod_{s\in S}\underbrace{\omega(\widehat{\vec A}_s)}_{\varepsilon(G_s^b)}
\]
\mspage{10}
